\chapter{Real ThermoDynamics as Model of Physics}\label{ch:essence}

\paragraph{What ThermoDynamics is.}
ThermoDynamics concerns dynamical transformations between kinetic
energy and heat energy in slightly viscous compressible fluids,
mediated by pressure work $W$ and turbulent dissipation $D$, and
realised by a computational physical process operating with finite
precision.

\paragraph{Turbulent dissipation.}
Turbulent dissipation transforms \emph{coordinated} bulk kinetic
energy at large scales into \emph{uncoordinated} small-scale kinetic
energy --- which, observed at the macroscopic level, is internal
heat energy. The transformation is not optional: it is the mechanism
by which the underlying continuum equations of compressible flow
avoid the development of infinite velocity gradients in the
convective instability characteristic of slightly viscous flow.
Without dissipation the flow would crash on its own resolution;
with dissipation it does not, and the price paid is that organised
motion at large scales is converted to disorganised motion at small
scales.

\paragraph{Partial irreversibility from precision deficit.}
The conversion is \emph{partially irreversible} for a precise
reason: to recreate coordinated motion from the resulting
uncoordinated small-scale motion would require a precision in the
small-scale degrees of freedom that the physical realisation does
not, and any finite-precision computation cannot, supply. The
irreversibility of the 2nd Law of Thermodynamics is the inevitability
of this precision deficit, no more and no less.

\paragraph{Physical computation, digital simulation.}
Physical ThermoDynamics is, in this reading, an analog computation
performed by Nature at finite precision --- the precision set by
whatever microscopic process limits the resolution at very short
length and time scales. The book mimics that physical computation
with a digital computation, also at finite precision --- the
precision set by the resolution of a mesh, the arithmetic of a
floating-point number, and the time step of an integrator. The two
computations agree on those output quantities that survive the
difference in precision: mean values, integrated fluxes, total
dissipation $D$, drag, lift, temperature gaps between chambers.
They disagree on point-values of turbulent flow, which neither
computation has the precision to predict.

\section{The Finite-Precision 2nd Law}\label{sec:fp2L}

The essence of the 2nd Law of Thermodynamics, on this reading, can be
stated in one sentence:

\begin{center}
\fbox{\parbox{0.86\textwidth}{\centering\itshape
The 2nd Law is the \textbf{necessary destruction of unstable
fine-scale structure} in the presence of \textbf{finite precision},
which \textbf{cannot be recovered} because the precision required to
reverse the destruction is not available.
}}
\end{center}

Each of the three emphasised clauses does load-bearing work.

\paragraph{(i) Necessary / unavoidable.}
The destruction is not an empirical observation about real fluids but
a structural consequence of finite resolution. A sufficiently sharp
velocity gradient is unstable under the smallest perturbations that
the available resolution can support, and no finite-precision dynamics
can prevent the instability from completing. The destruction
\emph{must} happen; nothing in the laws of motion can hold a structure
together below the resolution that supports it.

\paragraph{(ii) Instability of fine-scale structure (sharp gradients).}
The structure that gets destroyed is the inertially cascading
small-scale eddies of turbulent flow and the sharp fronts of shock
waves --- both characterised by velocity gradients that the underlying
continuum equations would, at infinite resolution, support
indefinitely but that finite resolution cannot represent.
Mathematically: the smooth-solution branch of the continuum equations
loses uniqueness past a critical scale; physically: the flow
\emph{breaks} into the small-scale motion the resolution permits.

\paragraph{(iii) Irreversibility from precision deficit.}
Once the bulk kinetic energy has been redistributed across the
resolution-limited small scales as incoherent thermal motion, the
precision required to coordinate the small-scale motion back into the
original bulk motion is not present in any finite-precision
representation --- neither the analog physical world nor the digital
simulation. The reverse process is not forbidden by the equations of
motion; it is forbidden by the precision the equations of motion are
solved with.

\bigskip

We refer to this formulation, throughout the book, as the
\textbf{Finite-Precision 2nd Law} --- \textbf{FP 2nd Law} for short
--- in distinction to the classical state-function 2nd Law of
equilibrium thermodynamics, the statistical-mechanical 2nd Law of
Boltzmann, and the information-theoretic 2nd Law of Shannon. The FP
2nd Law subsumes their content in the regimes where they apply and
extends to the regimes they do not. The remainder of this volume is
an extended demonstration of it: Joule's experiment, the
piston-cylinder, refrigeration and heat-engine cycles, chemical
bonding, blackbody radiation, the Big Bang and Big Crunch are all
instances of the same precision-deficit mechanism, played out in
different geometries and at different scales.
